%% notes-tex/database-expansion-100/sections/classes.tex
%% [[experiments.database.expansion-100]]
%% https://github.com/Mjvolk3/torchcell/tree/main/notes-tex/database-expansion-100/sections/classes.tex
%% SOURCE: hand-authored prose; row numbers and counts come from
%% experiments/database/scripts/build_candidate_datasets_table.py

\section{The classes\secstatus{todo}}\label{sec:classes}

Six classes were requested. Three more were added because the expansion turned up
coherent groups that did not fit any of them: regulatory DNA, deep mutational scans and
combinatorial genomes, all covered in Sec.~\ref{sec:depth}.

\subsection{Natural isolates, and where sequence diversity crosses over\secstatus{todo}}

Natural variation is the largest class after the backbone rows (Table~\ref{tab:counts}), and the
crossover the class was asked for is concentrated in three of them.

Dutta 2026 barcodes 520 natural isolates and profiles them against more than 600
compounds. Those isolates are drawn from the sequenced 1,011 panel, which is the same panel
behind the already-built Caudal 2024 pan-transcriptome and behind Muenzner 2024.
One genotype axis therefore carries a genome, a transcriptome, a proteome, and now a
chemogenomic response surface, which is the join that makes a transfer question askable at
all.

Hale 2024 is the sharpest instrument for that question. It crosses 8,046 CRISPRi
guides over 1,721 genes into 169 sequenced segregants, so it measures directly how much a
perturbation's effect depends on the background it is made in. Galardini 2019
asks the same question with deletions instead of knockdowns, building 3,786 knockouts
independently in four sequenced backgrounds across 38 conditions, and reports that about
18.5 percent of deletion phenotypes change between backgrounds.

That 18.5 percent is the reason this class matters beyond breadth. Nearly every fitness and
metabolite dataset in the built set is S288C, and a model trained only there has no way to
know which of its predictions survive a change of chassis.

\subsection{Tolerance on recalcitrant biomass, and product tolerance\secstatus{todo}}

Lignocellulosic inhibitor coverage in the built set currently runs through broad compound panels rather than dedicated screens, so
three rows close named gaps: Pereira 2014 screens the deletion collection against
real industrial wheat-straw hydrolysate rather than reconstituted single toxins, Endo 2008 covers vanillin, and Xiao 2014 covers furfural through a knockdown axis
that reaches essential genes a deletion collection cannot.

Product tolerance is answered mainly on isobutanol, and the earlier assumption that no
genome-wide isobutanol screen exists is wrong. Kuroda 2019 screens the deletion
collection for isobutanol-specific fitness against an ethanol comparator, so the
alcohol-general stress response subtracts out; its \gene{GLN3} hit raised production
4.9-fold, which is a rare case of a tolerance screen translating into a titer. Liu 2021 is an independent full-collection screen of the same phenotype with a different
readout, giving a same-genotype, same-phenotype, different-method replicate. Crook 2016 adds a dose-graded knockdown axis over isobutanol and 1-butanol together, and
Peeters 2021 supplies isobutanol titer as a mapped trait across 1,125 sequenced
segregants.

Van Leeuwen 2024 runs acetic, butyric and octanoic acid on one strain panel, which is what
lets shared and acid-specific mechanisms separate inside a single study instead of across
screens that cannot be compared.

\notesec{Acetic acid, production as well as tolerance} Acetic acid is now covered on both
sides, which is unusual. On tolerance: Mira 2010 screens the full deletion collection at
controlled pH and reports about 650 determinants, Mukherjee 2021 covers the essential genes
by CRISPRi, Sousa 2013 covers gain-of-function by overexpression, and the already-built
Mormino 2022 covers it by biosensor. Four perturbation types, one phenotype.

Production is thinner but real. Peeters 2021 measures acetic acid titer as one of its 18
mapped traits across 1,125 sequenced segregants, which is the only quantitative
genotype-to-acetate-titer map at scale. Chang 2012 screens the whole deletion collection for
altered organic-acid output including acetate, though the readout is an ordinal halo rather
than a titer. No genome-wide screen was found whose primary label is a quantitative acetate
titer per strain, so the pairing that exists is Peeters for production against the four
tolerance screens, not a matched production screen against them.

\subsection{CRISPR library screens\secstatus{todo}}

The already-built Lian 2019 MAGIC screen reads out furfural
tolerance by growth selection; Dong 2021 is the sibling that swaps the growth
selection for an S-adenosyl-methionine biosensor and sorts by fluorescence, so the label
becomes product concentration rather than fitness. That is the pattern worth copying for
every precursor in Sec.~\ref{sec:precursors}, and it is the biosensor MAGIC screen the
scoping request asked after.

Bao 2018 and Roy 2018 both move the genotype axis below the ORF, to
designed single-nucleotide and small-indel variants. Both are designed rather than
per-strain-verified, so both would enter as designed genotypes with a verification status.
McGlincy 2021 and Momen-Roknabadi 2020 are two independently designed
CRISPRi libraries; ingesting both makes library-to-library transfer a cheap and decisive
test of whether a model has learned gene function or guide-design idiosyncrasy.

\subsection{Expression across conditions\secstatus{todo}}

This class is genuinely thin in a way worth stating rather than papering over. Nadal-Ribelles 2025 is the largest genotype-indexed expression dataset in
yeast at roughly 3,500 knockouts under osmotic stress, and it is already built. The
candidates here supply what it lacks rather than more of it: Gasch 2000 is the
canonical wild-type stress-response baseline that any genotype-indexed transcriptome has to
be read against, N'Guessan 2025 profiles roughly 4,500 sequenced segregants by
single-cell RNA-seq, and Leutert 2023 covers 101 conditions at the phosphosite
level, where enzyme regulation acts faster than transcription.

The initial reading was that datasets in the microbe-perturb-seq collection carry little
data over large genotype numbers. That holds. Nothing found in this pass changes it, and
the gap it points to is addressed in Sec.~\ref{sec:perturbseq} rather than by another row.

\subsection{Metabolite and precursor coverage\secstatus{todo}}\label{sec:precursors}

The precursor accounting in
\file{notes/metabolism.central-carbon-precursors.md} sets the target: of the twelve central
carbon precursors, acetyl-CoA is measured on 19 strains and succinyl-CoA and malonyl-CoA on
none, so identity coverage is good and statistical power is not.

Blank 2005 is the only dataset in this table that measures flux rather than a
concentration or a titer, across roughly 200 deletion mutants, and it is the only real
handle on acetyl-CoA turnover under genetic perturbation. Malonyl-CoA still has no direct
measurement anywhere; the nearest handles remain its dominant sinks, which the already-built
Xue 2025 free fatty acids and da Silveira 2014 lipidomics supply, and Zhu 2014
extends with a lipidome across 129 kinase and phosphatase deletions.

Cooper 2010 is the highest-value row in the class for a reason unrelated to
novelty. It measures the same amino-acid pools as the already-built Mulleder 2016, on an
overlapping genotype axis, by capillary electrophoresis instead of mass spectrometry. Two
independent platforms measuring one trait class is the cleanest available test of whether a
model has learned biology or a batch effect, and nothing else in the built set offers it.

On carotenoids, Ozaydin 2013 is built. Trikka 2015 is the other high-throughput
carotenoid screen: 4,700 heterozygous deletions scored by color, then carried through
iterative rounds to a sclareol titer. Halving gene dosage rather than deleting outright is
what finds the isoprenoid flux control points a null would kill, and no other row in this
table uses that design. It is currently recorded as figure-only, so confirming that
Additional file 1 carries per-strain scores is a precondition, not a formality.

\subsection{Regulatory DNA, deep scans and combinatorial genomes\secstatus{todo}}

Three new classes, all with exactly known sequences and all going deep on
a locus rather than shallow across a genome. They are covered in Sec.~\ref{sec:depth},
because their value is best stated against the Perturb-seq specification rather than as
another phenotype class.

\subsection{Perturbation modality and reference backbone\secstatus{todo}}

None of these is a headline dataset and one is load-bearing.
Puddu 2019 is whole-genome sequence for every strain in the deletion collection.
Twenty-one recommended rows declare an \term{S288C-KO} sequence basis, and without Puddu
that basis is an assumption that each strain differs from the reference at exactly the
intended locus. Puddu shows it frequently does not. Ingesting it converts a shared
assumption across half this table into a measurement, which is worth more than its own
4,800 instances suggest.
